% 第一章 緒論
\chapter{緒論}\label{ch:intro}

\section{研究背景：功耗成為算力的上限}\label{sec:motivation-cap}

大型計算系統的電力與散熱都有預算。機房依機櫃的供電能力設定功耗上限，學生超算競賽（Student Cluster Competition, SCC）限制整個隊伍的總功率，NVIDIA GH200 superchip 也允許使用者直接設定整顆模組（GPU、CPU 與記憶體）的功耗上限（本叢集的驅動程式接受 200–1000~W，見第 \ref{sec:platform} 節）。當程式的耗電碰到上限時，GPU 的功耗管理會降低時脈與電壓，直到功耗不超過上限。因此在功耗受限時，效能不再由硬體的最高時脈決定，而是由每次運算所消耗的能量決定：每次運算的耗電越低，時脈便能拉得越高。

\section{功耗隨資料而變}\label{sec:motivation-data}

CMOS 電路的動態功耗約為
\begin{equation}
  P_{\mathrm{dyn}} \approx a\, C\, V^{2} f ,
  \label{eq:cmos}
\end{equation}
其中 $C$ 為負載電容、$V$ 為電壓、$f$ 為時脈，$a$ 為切換活動率（本報告的 $\alpha$ 另指第 \ref{sec:errmodel} 節的誤差係數）：資料位元在相鄰時脈之間翻轉（0 變 1 或 1 變 0）得越多，電路充放電的次數越多，耗電就越高。因此同一個程式、同樣的運算量，輸入資料不同，耗電也不同。前人已在 A100 與 GH200 上量到，同樣的矩陣乘法，隨機輸入比固定值或全零輸入多耗 60–67\% 的功率（見第三章）。

\section{應用用不到全部的尾數}\label{sec:motivation-bits}

FP64 有 52 個尾數位元，約相當於 16 位十進位有效數字，但許多應用並不需要這麼高的精度：HPL 只要求殘差小於門檻；平面波密度泛函計算的自洽迭代在前期只需要很寬鬆的精度；AI 推論甚至以 FP8（只有 3 個尾數位元）也能維持準確度。這些用不到的低位元在運算中仍持續隨機翻轉，所消耗的電力並未帶來有用的精度。

\section{研究問題}\label{sec:questions}

綜合以上三點，本專題的核心想法是：\textbf{以軟體將用不到的低位尾數設為 0，使其不再翻轉，並在功耗上限下把省下的功耗轉換為效能。}本專題要回答以下三個問題：

\begin{enumerate}
  \item \textbf{機制：}為什麼少幾個尾數位元就會省電？省電來自位元不再翻轉，還是來自數值為 0？
  \item \textbf{效益：}在功耗上限下，真實的科學應用能換得多少效能與省電？需要付出多少精度？
  \item \textbf{適用範圍：}與直接改用較小格式的做法（INT8 模擬 FP64、FP32／3xTF32、FP8）相比，尾數歸零在哪些情況下有價值？
\end{enumerate}

\section{主要貢獻}\label{sec:contrib}

\begin{enumerate}
  \item \textbf{功耗上限下真實應用的收益。}以「未修改的程式需要多給多少功耗上限才能追上」的等效功耗量化收益，涵蓋附信賴區間的矩陣乘法微基準、調參後的 8 節點 HPL、HPL-MxP、Quantum ESPRESSO、高功耗的 DF-MP2，以及指數跨度寬且功耗高的高斯過程地統計，並以各程式本身的判準（殘差、總能量、估計參數）檢查結果。
  \item \textbf{省電機制。}以全 1 或固定樣式取代 0 填入被捨去的位元，在 tensor core 矩陣乘法、FP64 乘加單元與記憶體搬移三個層次上，分開「位元不再翻轉」與「數值為 0」兩種效應，在 GH200 的模組功耗上量化了前人提出的翻轉假說。
  \item \textbf{所需位數與失效邊界。}提出只用求解器自身容許誤差決定位元數的規則，使 Quantum ESPRESSO 的總能量印出的每一位都不變；並測試 HPL（多組問題大小）、PySCF、第三方非矩陣乘法程式、LoRA 微調與六個語言模型可以截斷到什麼程度，以及捨入方式的影響。
  \item \textbf{與換格式做法的分界。}與 INT8 模擬（cuBLAS、GEMMul8）、FP32／3xTF32 及 AI 的 FP8 比較，指出分界點隨矩陣形狀、功耗上限、累加精度與指數跨度移動；列內指數跨度寬時，固定位數的 INT8 模擬可能安靜地算錯，而這個風險可以事先以低成本的逐列跨度統計偵測。
\end{enumerate}

\section{報告架構}\label{sec:outline}

第二章介紹浮點格式、GPU 功耗與功耗上限等背景知識；第三章整理相關研究；第四章說明研究方法與實驗平台；第五章呈現實驗結果，包括高功耗的 DF-MP2、指數跨度實驗，以及跨度寬且功耗高的高斯過程與 ExaGeoStat；第六章討論適用範圍與限制；第七章為結論與未來工作。附錄 A 列出實驗環境、程式位置與重現方式。
